\documentclass[../main]{subfiles}
\begin{document}
\chapter*{Planificación \materia}
\section{Diagnóstico 2023}
El diagnóstico áulico nos permite conocer cuáles son los ritmos y estilos de aprendizajes, nivel de dominio de las capacidades, los intereses de los estudiantes, conocer sus habilidades y debilidades y a partir de allí planificar las estrategias adecuadas al grupo clase. 
\info{\materia}{
Espacio Curricular: \materia  

Año: \ano

División: \division  

Profesor/a: \nombre
}

\subsection{DIAGNÓSTICO:} 

Marcar las capacidades a evaluar:  


\begin{itemize}
  \item[$\square$] Observación 
  \item[$\square$] Identificación 
  \item[$\square$] Comparación   
  \item[$\square$] Análisis
  \item[$\square$] Resumen
  \item[$\square$] Síntesis
  \item[$\square$] Clasificación
  \item[$\square$] Interpretación
  \item[$\square$] Formulación de críticas
  \item[$\square$] Búsqueda de suposiciones
  \item[$\square$] Imaginación y creatividad 
  \item[$\square$] Organización de datos  
  \item[$\square$] Formulación de hipótesis  
\end{itemize}
  

 

Explicitar los aprendizajes específicos evaluados: 

 

Instrumento/s utilizado/s:  

 

Descripción de los destinatarios: La característica del grupo clase frente a la tarea: fortalezas y debilidades en el período de diagnóstico. 

 

Alumnos que presentan alguna dificultad 

 \begin{table}[h]
     \centering
     \begin{tabular}{|c|c|}
         Alumnos & Dificultad \\
          & 
     \end{tabular}
   \end{table}


\section{PLANIFICACIÓN 2023}

ÁREA: TÉCNICA  

ESPACIO CURRICULAR:  

CURSO: $5^o 2^2$ 

DOCENTE: Ferreyra Gustavo David

ORIENTACIÓN EN: ELECTROMECÁNICA

ACUERDOS INSTITUCIONALES: 
ACUERDOS DEL ÁREA: 

\begin{table}[h]
    \centering
    \begin{longtable}{|c|c|c|c|}
        EJE & SABERES & ESTRATEGIAS & EVALUACIÓN\\
         & & & \\
    \end{longtable}
\end{table}


MARCAR CON NEGRITA LOS SABERES INVOLUCRADOS EN LOS PROYECTOS (*) 

 

Ajuste-Observaciones 1er periodo: 

Ajuste-Observaciones 2do periodo: 
\end{document}

